Carl Jung, a contemporary, broke with Freud but remained within the psychoanalytic framework, viewing dreams as direct mental expressions, compensations for waking life, and expressions of a "collective unconscious" through mythic archetypes \citep{jung1974dreams}. Until the 1950s, dream research was entirely subjective and interpretive, operating within this psychoanalytic paradigm.

The single most important event in modern sleep science occurred in 1953. Researchers \citet{aserinsky1953rem} discovered that dreaming "typically occur[s] during rapid eye movement (REM) sleep". This discovery was a paradigm shift.

For the first time, the subjective, internal state of dreaming was linked to an \textit{objective, physiological correlate} that could be measured in a laboratory. This finding moved oneirology (the study of dreams) from the analyst's couch to the neuroscience lab, providing a "foundation" for all subsequent scientific investigation.